\chapter{Conclusiones}
\label{sec:conclusiones}

El desarrollo y verificación del firmware sobre el microcontrolador STM32F103 (núcleo ARM Cortex-M3) junto con el
sistema operativo FreeRTOS permitieron cumplir satisfactoriamente con los requerimientos operativos planteados para el
Hackathon 2026.
La división funcional en tres tareas concurrentes con prioridades diferenciadas demostró ser eficaz para desacoplar
las responsabilidades críticas de control temporal frente a los procesos de diagnóstico analógico y comunicación
serie.

En particular, la estrategia de arranque seguro con retardo de despeje en Todo Rojo y el submuestreo temporal en
bloques de $100\ms$ dentro de la máquina de estados garantizaron una respuesta inmediata ante cambios de modo o
consignas recibidas por la red serie, evitando estados intermedios inconsistentes o demoras perjudiciales en la
señalización vial. Asimismo, la incorporación del algoritmo de cálculo de corriente esperada posibilitó la detección
precoz de anomalías por lámparas abiertas o sobrecorrientes en los cabezales, confiriendo al equipo capacidades de
diagnóstico preventivo.

Por su parte, la implementación del protocolo estándar Modbus RTU sobre la interfaz USART1, asistida por una cola de
interrupción y detección de fin de trama por silencio de línea ($t_{3.5}$), aseguró una comunicación confiable con el
panel de supervisión web sin sobrecargar la CPU. En definitiva, la solución técnica alcanzada combina robustez en la
conmutación física de fases, flexibilidad de configuración remota y supervisión activa del estado del sistema.
